\section{Quantum loop algebras and their finite-dimensional representations}
\label{Sec : qloop}

\subsection{Quantum loop algebras}
\label{ssec : quantum loop algebras}

Let $\mathfrak{g}$ be a complex finite-dimensional simple Lie algebra
with the set of Dynkin indices $I = \{1, \ldots, n\}$
and Cartan matrix $C = (c_{ij})_{i, j \in I}$.
Let $r$ be the maximum number of edges in the Dynkin diagram
and $D = \mathrm{diag}(d_i \mid i \in I)$ be the unique diagonal matrix
such that $d_i \in \{1, r\}$ for each $i \in I$
and $DC = (d_i c_{ij})_{i, j \in I}$ is symmetric.

Let $q$ be an indeterminate and $\kk$
the algebraic closure of the rational function field $\Q(q)$.
Let $q_i \coloneq q^{d_i}$.
For $k \in \Z$, we set 
\[
[k]_{q_i} \coloneq \frac{q_i^k - q_i^{-k}}{q_i - q_i^{-1}}. 
\]
For integers $k \geq l \geq 0$, we set
\[
[k]_{q_i}! \coloneq [k]_{q_i}[k - 1]_{q_i}\cdots [2]_{q_i}[1]_{q_i},
\quad
\left[\begin{matrix} k \\ l \end{matrix}\right]_{q_i}
\coloneq \frac{[k]_{q_i}!}{[k - l]_{q_i}! [l]_{q_i}!}.
\]

\begin{Def}
\label{Def : quantum loop algebra}
\emph{The quantum loop algebra} $U_q(L\mathfrak{g})$
is the $\kk$-algebra given by the following generators and relations :

generators : \quad
$x_{i, k}^{\pm}\ (i \in I, k \in \Z), \quad
k_i^{\pm 1}\ (i \in I), \quad
h_{i, l}\ (i \in I, l \in \Z \setminus \{0\})$

relations :
\begin{itemize}
    \item $k_i k_i^{-1} = k_i^{-1} k_i = 1,
    k_i k_j = k_j k_i$ for $i, j \in I$,
    \item $k_i x_{j, k}^{\pm} k_i^{-1}
    = q_i^{\pm c_{ij}} x_{j, k}^{\pm}$
    for $i, j \in I$ and $k \in \Z$,
    \item $[k_i, h_{j, l}] = [h_{i, l}, h_{j, m}] = 0$
    for $i, j \in I$ and $l, m \in \Z \setminus \{0\}$,
    \item $[h_{i, l}, x_{j, m}^{\pm}]
    = \pm \frac{1}{l}[l c_{ij}]_{q_i}x_{j, l + m}^{\pm}$
    for $i, j \in I, l \in \Z \setminus \{0\}$ and $m \in \Z$,
    \item $x_{i, k + 1}^{\pm} x_{j, l}^{\pm}
    - q_i^{\pm c_{ij}} x_{j, l}^{\pm} x_{i, k + 1}^{\pm}
    = q_i^{\pm c_{ij}} x_{i, k}^{\pm} x_{j, l + 1}^{\pm}
    - x_{j, l + 1}^{\pm} x_{i, k}^{\pm}$
    for $i, j \in I$ and $k, l \in \Z$,
    \item $[x_{i, k}^{+}, x_{j, l}^{-}]
    = \displaystyle\delta_{i, j}
    \frac{\phi_{i, k + l}^{+} - \phi_{i, k + l}^{-}}{q_i - q_i^{-1}}$
    for $i, j \in I$ and $k, l \in \Z$,
    \item $\displaystyle\sum_{\sigma \in \mathfrak{S}_b}\sum_{k = 0}^{b}
    (-1)^k \left[\begin{matrix} b \\ k \end{matrix}\right]_{q_i}
    x_{i, r_{\sigma(1)}}^{\pm}\cdots x_{i, r_{\sigma(k)}}^{\pm}
    x_{j, s}^{\pm} x_{i, r_{\sigma(k + 1)}}^{\pm}\cdots
    x_{i, r_{\sigma(b)}}^{\pm} = 0$,
    where $b = 1 - c_{ij}$, for $i, j \in I$ with $i \neq j$
    and $s, r_1, \ldots, r_b \in \Z$,
\end{itemize}
where
\[
\phi_i^{\pm}(z) \coloneq
\sum_{r = 0}^{\infty}\phi_{i, \pm r}^{\pm}z^{\pm r} \coloneq
k_i^{\pm 1}\exp\left(\pm (q_i - q_i^{-1})
\sum_{l = 1}^{\infty}h_{i, \pm l}z^{\pm l}\right).
\]
\end{Def}

\subsection{Finite-dimensional representations}
\label{ssec : finite-dimensional representations}

A finite-dimensional $U_q(L\mathfrak{g})$-module is said to be of type
$\mathbf{1}$
if the element $k_i$ acts as a diagonal $\kk$-linear operator
whose eigenvalues belong to the set $\{q^k \mid k \in \Z\}$
for each $i \in I$.
It is well-known that the study of finite-dimensional representations
of $U_q(L\mathfrak{g})$ reduces essentially to the study of the category
$\mathscr{C}_{\mathfrak{g}}$ of type $\mathbf{1}$ finite-dimensional
$U_q(L\mathfrak{g})$-modules.
It is also well known that $U_q(L\mathfrak{g})$ is a subquotient of the
corresponding
quantum affine algebra and the study of type $\mathbf{1}$
finite-dimensional modules of quantum affine algebra reduces
to the study of $\mathscr{C}_{\mathfrak{g}}$.
Since $U_q(L\mathfrak{g})$ has a Hopf algebra structure,
$\mathscr{C}_{\mathfrak{g}}$ is a monoidal category.

Let $V \in \mathscr{C}_{\mathfrak{g}}$.
Since the elements
$\{k_i^{\pm 1}, h_{i, l} \mid i \in I, l \in \Z \setminus \{0\}\}$
mutually commute, there is a decomposition
\[
V = \bigoplus_{\gamma \in
(\kk[\![z]\!] \times \kk[\![z^{-1}]\!])^I}V_{\gamma}
\]
where, for each
$\gamma = (\gamma_i^{+}(z), \gamma_i^{-}(z))_{i \in I}$,
we define $V_{\gamma}$ as the subspace of $V$ on which each coefficient
of the series $\phi_i^{\pm}(z) - \gamma_i^{\pm}(z)$ acts nilpotently
for any $i \in I$.
If $V_{\gamma} \neq 0$, it is called \emph{an $\ell$-weight space}
and $\gamma$ is called \emph{$\ell$-weight}.

For a simple module $L$ in $\mathscr{C}_{\mathfrak{g}}$,
there uniquely exists an $\ell$-weight space $L_{\gamma_0}$
such that $x_{i, k}^{+}L_{\gamma_0} = 0$
for all $i \in I$ and $k \in \Z$.
Such $\ell$-weight $\gamma_0$ is called
\emph{the highest $\ell$-weight} of $L$. 
The isomorphism classes of simple objects
in $\mathscr{C}_{\mathfrak{g}}$ is classified by highest $\ell$-weights :

\begin{Thm}[{\cite[Theorem 3.3]{CP95}, \cite[Theorem 12.2.6]{CP94}}]
\label{Thm : Drinfeld poly}
If $\gamma = (\gamma_i^{+}(z), \gamma_i^{-}(z))_{i \in I}$
is an highest $\ell$-weight of a simple $U_q(L\mathfrak{g})$-module
in $\mathscr{C}_{\mathfrak{g}}$, for each $i \in I$ there exists a unique
polynomial
$P_i(z) \in 1 + z\kk[z]$ for $i \in I$ such that
\begin{equation}
\label{eq : ellhwt}
\gamma_i^{\pm}(z) = q_i^{\deg(P_i)}
\frac{P_i(zq_i^{-1})}{P_i(zq_i)}.
\end{equation}
Conversely, for any $(P_i(z))_{i \in I} \in (1 + z\kk[z])^I$,
we have a simple $U_q(L\mathfrak{g})$-module in $\mathscr{C}_{\mathfrak{g}}$
whose highest $\ell$-weight $\gamma$ is given by
{\rm (\ref{eq : ellhwt})}.
\end{Thm}

This $I$-tuple of polynomials $(P_i(z))_{i \in I}$
associated with a simple module $L$
is called \emph{the Drinfeld polynomials of $L$}.


Let $\Irr\mathscr{C}_{\mathfrak{g}}$ denote the classes of simple objects of
$\mathscr{C}_{\mathfrak{g}}$.
By theorem~\ref{Thm : Drinfeld poly}, $\Irr\mathscr{C}_{\mathfrak{g}}$
is parametrized by $(1 + z\kk[z])^I$.

Let $Y_{i, a}\ (i \in I, a \in \kk^{\times})$ be an indeterminate and
\[
\mathcal{M}_{+} \coloneq
\biggl\{m = \prod_{i \in I, a \in \kk^{\times}}Y_{i, a}^{u_{i, a}}
\ \bigg|\ u_{i, a} \in \Z_{\geq 0},
u_{i, a} = 0\ \text{except for finitely many }(i, a)\biggr\}.
\]
Then, there is a bijection
\[
\mathcal{M}_{+} \to (1 + z\kk[z])^I; \quad
\prod_{i \in I, a \in \kk^{\times}}Y_{i, a}^{u_{i, a}}
\mapsto \left(\prod_{a \in \kk^{\times}}(1 - az)^{u_{i, a}}\right)_{i \in I}.
\]
Using this bijection, $\Irr\mathscr{C}_{\mathfrak{g}}$ is also parametrized
by $\mathcal{M}_{+}$.
Let $L(m)$ denote the simple module corresponding to the monomial $m$.
The monomial $m$ is called the \emph{highest $\ell$-weight monomial}
of the simple module $L(m)$.

\begin{Prop}[{\cite[Corollary 3.5]{CP95}}]
\label{Prop : subquot}
Let $m_1, \ldots, m_k \in \mathcal{M}_{+}$.
Then, $L(m_1\cdots m_k)$ is isomorphic to a subquotient
of $\bigotimes_{i = 1}^{k}L(m_i)$.
\end{Prop}

For each $a \in \kk^{\times}$, we consider a $\kk$-algebra automorphism
$\tau_a$ of $U_q(L_{\mathfrak{g}})$ given by
\[
\tau_a(k_i) = k_i, \quad
\tau_a(x_{i, k}^{\pm}) = a^k x_{i, k}^{\pm}, \quad
\tau_a(h_{i, l}) = a^l h_{i, l},
\]
for $i \in I, k \in \Z, l \in \Z \setminus \{0\}$.
This automorphism $\tau_a$ is a $q$-analog of the loop rotation
$z \mapsto az$ for the loop algebra $L\mathfrak{g} = \mathfrak{g}
\otimes_{\C}\C[z^{\pm}]$.
For a finite-dimensional representation $V$ in $\mathscr{C}_{\mathfrak{g}}$,
the twisted representation $\tau_a^*V$ by $\tau_a$
again belongs to $\mathscr{C}$.
Thus the assignment $V \mapsto \tau_a^*V$ defines a monoidal
auto-equivalence $T_a$ of $\mathscr{C}_{\mathfrak{g}}$.
If we denote by $m \mapsto m_a$ the monoid automorphism
of $\mathcal{M}_{+}$ given by $Y_{i, b} \mapsto Y_{i, ab}$,
we have $T_a(L(m)) \cong L(m_a)$
for any simple module $L(m) \in \mathscr{C}_{\mathfrak{g}}$.

In $\mathscr{C}_{\mathfrak{g}}$,
$V \otimes W$ is not necessarily isomorphic to $W \otimes V$
for $V, W \in \mathscr{C}_{\mathfrak{g}}$.
However, using the $q$-character,
which is a type of character in the quantum loop algebra,
we have the following fact.

\begin{Prop}[{\cite[Theorem 3]{FR98}}]
\label{Prop : q-character}
Let $K(\mathscr{C}_{\mathfrak{g}})$ be the Grothendieck ring
of the monoidal category $\mathscr{C}_{\mathfrak{g}}$.
Then, there is an injective algebra homomorphism
$\chi_q \colon K(\mathscr{C}_{\mathfrak{g}}) \to
\Z[Y_{i, a}^{\pm 1} \mid i \in I, a \in \kk^{\times}]$.
In particular, although $\mathscr{C}_{\mathfrak{g}}$ is not braided,
$K(\mathscr{C}_{\mathfrak{g}})$ is commutative.
\end{Prop}

\subsection{Some subcategories of $\mathscr{C}_{\mathfrak{g}}$}
\label{ssec : subcategories}

For $i, j \in I$, write
\[
d_{ij} \coloneq \min\{d_i, d_j\}.
\]
A function $\epsilon \colon I \to \{0, 1\}$
is called \emph{a parity function} for $\mathfrak{g}$
if it satisfies the condition
\[
\epsilon_i \equiv \epsilon_j + d_{ij} \pmod{2}
\]
for any adjacent nodes $i, j \in I$.
In what follows, we choose and fix a parity function $\epsilon$.
Let $\widehat{I} \coloneq
\{(i, p) \in I \times \Z \mid p - \epsilon_i \in 2\Z\}$.

\begin{Def}[{\cite[$\S$3.7]{HL10}, \cite[Definition 2.6]{FHOO22}}]
\label{Def : Cg0}
Let $\mathscr{C}_{\mathfrak{g}}^{0}$ be the smallest full subcategory of
$\mathscr{C}_{\mathfrak{g}}$
containing $L(Y_{i, q^p})$ for $(i, p) \in \widehat{I}$
and trivial module, and stable under taking subquotients,
extensions, and tensor products.
\end{Def}

\begin{Prop}[{\cite[$\S$3.7]{HL10}, \cite[Proposition 2.8]{FHOO22}}]
\label{Prop : Cg factorization}
Let $L \in \Irr\mathscr{C}_{\mathfrak{g}}$.
Then $L$ is factorized into a commutative tensor product
\[
L \cong T_{a_1}(L_1) \otimes \cdots \otimes T_{a_d}(L_d)
\]
with $L_k \in \Irr\mathscr{C}_{\mathfrak{g}}^{0}$ such that $L_k$ is non-trivial and $a_k \in \kk^{\times}$
for $1 \leq k \leq d$ such that
$a_k/a_l \notin q^{2\Z}$ for all $k \neq l$.
Moreover, the set $\{T_{a_k}(L_k)\}_k$ is unique. 
\end{Prop}

For a complex finite-dimensional simple Lie algebra $\mathfrak{g}$,
we associate a Lie algebra $\mathfrak{g}_{\mathrm{fin}}$
of simply laced type as follows :
\[
\begin{array}{c||c|c|c|c|c|c|c}
\mathfrak{g} & A_n\ (n \geq 1) & B_n\ (n \geq 2) & C_n\ (n \geq 3)
& D_n\ (n \geq 4) & E_{6, 7, 8} & F_4 & G_2 \\
\hline
\mathfrak{g}_{\mathrm{fin}} & A_n & A_{2n - 1} & D_{n + 1}
& D_n & E_{6, 7, 8} & E_6 & D_4
\end{array}
\]
Let $\Delta_{\mathfrak{g}_{\mathrm{fin}}}$ be the Dynkin diagram
for $\mathfrak{g}_{\mathrm{fin}}$.
Then there exists a Dynkin diagram automorphism $\sigma$
of $\Delta_{\mathfrak{g}_{\mathrm{fin}}}$ whose orbit set
$\Delta_{\mathfrak{g}_{\mathrm{fin}}}^{\sigma}$
yields the Dynkin diagram $\Delta_{\mathfrak{g}}$.
We denote by $I_{\mathrm{fin}} = \{1, 2, \ldots, |I_{\mathrm{fin}}|\}$
the index set of $\Delta_{\mathfrak{g}_{\mathrm{fin}}}$.
To distinguish the index set $I_{\mathrm{fin}}$ and $I$,
we use $\iota$ for indices in $I_{\mathrm{fin}}$
and $i$ for indices in $I$.
Let $\pi \colon I_{\mathrm{fin}} \to I$
be the projection induced by $\sigma$.
For $\iota \in I_{\mathrm{fin}}$,
we write $d_{\iota} \coloneq |\pi^{-1}\pi(\iota)|$.

\begin{Def}[{\cite[Definition 6.9]{KKOP24}}]
\label{Def : admissible sequence}
We say that an infinite sequence
$\mathfrak{s} \coloneq ((\iota_k, p_k))_{k \in \Z}$
in $I_{\mathrm{fin}} \times \Z$ is \emph{admissible}
if the following conditions are satisfied :
\begin{enumerate}
    \item[(1)] $p_{s^{+}} = p_s + 2d_{\iota_s}$,
    \item[(2)] $p_t = p_s + \min\{d_{\iota_s}, d_{\iota_t}\}$
    if $\iota_s$ and $\iota_t$ are adjacent
    and $t^{-} < s < t < s^{+}$,
    \item[(3)] $p_t - (p_s + 2) \in 2d_{\iota_s}\Z$
    if $s, t \in \Z$ satisfy $\iota_t = \sigma(\iota_s)$,
    \item[(4)] $s_{\iota_k}s_{\iota_{k + 1}}\cdots
    s_{\iota_{k + l - 1}} = w_0$ for every $k \in \Z$,
    where $l$ denotes the length of the longest element $w_0$
    of the Weyl group of $\mathfrak{g}_{\mathrm{fin}}$.
\end{enumerate}
Here, for $s \in \Z$,
$s^{+} \coloneq \min\{t > s \mid \iota_t = \iota_s\}$
and $s^{-} \coloneq \max\{t < s \mid \iota_t = \iota_s\}$.
\end{Def}

In what follows, we only consider admissible sequences
$\mathfrak{s} = ((\iota_k, p_k))_{k \in \Z}$
which satisfies $p_k \equiv \epsilon_{\pi(\iota_k)} \pmod{2}$.

\begin{Def}[{\cite[Definition 6.16]{KKOP24}}]
\label{Def : category}
Let $[a, b]$ be an interval with $-\infty \leq a \leq b \leq \infty$.
Let $\mathfrak{s} = ((\iota_k, p_k))_{k \in \Z}
\subset I_{\mathrm{fin}} \times \Z$ be an admissible sequence.
We define the category $\mathscr{C}_{\mathfrak{g}}^{[a, b], \mathfrak{s}}$
as the smallest full subcategory of $\mathscr{C}_{\mathfrak{g}}^{0}$ which
contains
$S_k^{\mathfrak{s}} \coloneq L(Y_{\pi(\iota_k), q^{p_k}})$
$(k \in [a, b])$ and trivial module,
and is closed under taking subquotients, extensions, and tensor products.
\end{Def}
We call $S_k$ a \emph{cuspidal module}. 

\begin{Prop}[{\cite[$\S$6.3]{KKOP24}}]
\label{Prop : whole category}
For any admissible sequence $\mathfrak{s}$,
$\mathscr{C}_{\mathfrak{g}}^{[-\infty, \infty], \mathfrak{s}} =
\mathscr{C}_{\mathfrak{g}}^{0}$.
\end{Prop}

\begin{Prop}[{\cite[$\S$6.3]{PBW24}}]
\label{Prop : PBW parameterization}
Let $[a, b]$ be an interval and
$\mathfrak{s} = ((\iota_k, p_k))_{k \in \Z}$
be an admissible sequence.
\begin{enumerate}
    \item[(i)] For any sequence $(c_k)_k \in \Z_{\geq 0}^{[a, b] \cap \Z}$
    with $c_k = 0$ for all but finitely many $k$'s,
    the head of $S_b^{c_b} \otimes S_{b - 1}^{c_{b - 1}}
    \otimes \cdots \otimes S_a^{c_a}$
    is in $\Irr\mathscr{C}_{\mathfrak{g}}^{[a, b], \mathfrak{s}}$.
    Moreover, this correspondence form a bijection between
    \[
    \bigl\{(c_k)_k \in \Z_{\geq 0}^{[a, b] \cap \Z}
    \mid c_k = 0\ \text{for all but finitely many }k\text{'s}\bigr\}
    \]
    and $\Irr\mathscr{C}_{\mathfrak{g}}^{[a, b], \mathfrak{s}}$.
    \item[(ii)] The Grothendieck ring
    $K(\mathscr{C}_{\mathfrak{g}}^{[a, b], \mathfrak{s}})$ is the polynomial ring
    generated by $\{[S_s] \mid s \in [a, b]\}$
    and the set of highest $\ell$-weights
    of $\Irr\mathscr{C}_{\mathfrak{g}}^{[a, b], \mathfrak{s}}$ is the set of
    monomials
    in $Y_{\pi(\iota_k), q^{p_k}}$ $(k \in [a, b])$.
\end{enumerate}
\end{Prop}

For $L \in \Irr\mathscr{C}_{\mathfrak{g}}^{[a, b], \mathfrak{s}}$,
the corresponding sequence $(c_k)_k$ is called \emph{PBW parameter}
and we write it as $\mathrm{PBW}(L)$.
